\documentclass[11pt]{article}

% Shared notes settings for Elements of AIML lectures (UPES).
% Keep this file free of \documentclass to allow reuse via \input.

\usepackage[utf8]{inputenc}
\usepackage[T1]{fontenc}
\usepackage[a4paper,margin=1in]{geometry}
\usepackage{hyperref}
\usepackage{xurl}
\usepackage{booktabs}
\usepackage{amsmath}
\usepackage{amssymb}
\usepackage{listings}
\usepackage{xcolor}
\usepackage{enumitem}
\usepackage{parskip}

% Code listings (general-purpose).
\lstset{
  basicstyle=\ttfamily\small,
  keywordstyle=\color{blue},
  commentstyle=\color{gray},
  stringstyle=\color{teal},
  showstringspaces=false,
  columns=fullflexible,
  frame=single,
  framerule=0.2pt,
  breaklines=true
}

\setlist[itemize]{topsep=0.3em,itemsep=0.2em}
\setlist[enumerate]{topsep=0.3em,itemsep=0.2em}

% Simple per-section source line.
\newcommand{\notesources}[1]{\par\noindent{\small\color{gray}\textit{Sources:} \sloppy #1\par}}

% Unit-wise source strings for this subject (used by slides/notes).
% Path is relative to the lecture `latex/` folder (the compilation CWD).
\IfFileExists{../../../_shared/aiml-sources.tex}{%
  % Source strings for Elements of AIML (used by slides + notes).

\newcommand{\AIMLRepoURL}{https://github.com/tali7c/Elements-of-AIML}

\newcommand{\AIMLLabSources}{%
Provost and Fawcett, \textit{Data Science for Business}.;%
McKinney, \textit{Python for Data Analysis}.;%
WEKA Documentation (Waikato Environment for Knowledge Analysis), accessed 2026-02-28.;%
Matplotlib and Seaborn documentation, accessed 2026-02-28.%
}

\newcommand{\AIMLUnitISources}{%
Rich and Knight, \textit{Artificial Intelligence}, McGraw Hill, 2017.;%
Russell and Norvig, \textit{Artificial Intelligence: A Modern Approach} (4e), Ch. 1--2 (Introduction, Intelligent Agents).;%
Poole and Mackworth, \textit{Artificial Intelligence: Foundations of Computational Agents} (2e), selected sections.%
}

\newcommand{\AIMLUnitIISources}{%
Russell and Norvig, \textit{Artificial Intelligence: A Modern Approach} (4e), Ch. 7--9 (Logical Agents, First-Order Logic, Inference).;%
Huth and Ryan, \textit{Logic in Computer Science} (2e), selected sections on propositional and first-order logic.;%
Genesereth and Nilsson, \textit{Logical Foundations of Artificial Intelligence}, selected chapters.%
}

\newcommand{\AIMLUnitIIISources}{%
Mueller and Massaron, \textit{Machine Learning for Dummies}, For Dummies, 2016.;%
Mitchell, \textit{Machine Learning}, selected chapters on learning basics and evaluation.;%
Scikit-learn documentation, accessed 2026-02-28.%
}

\newcommand{\AIMLUnitIVSources}{%
Mueller and Massaron, \textit{Machine Learning for Dummies}, For Dummies, 2016.;%
James, Witten, Hastie, and Tibshirani, \textit{An Introduction to Statistical Learning} (2e), selected chapters.;%
Scikit-learn documentation, accessed 2026-02-28.%
}

\newcommand{\AIMLUnitVSources}{%
NIST, \textit{AI Risk Management Framework (AI RMF 1.0)}, 2023.;%
Barocas, Hardt, and Narayanan, \textit{Fairness and Machine Learning} (online book), selected sections.;%
Knaflic, \textit{Storytelling with Data}, selected chapters on communicating results.%
}
%
}{%
  % Faculty copies live under `private/` so their compile CWD is different.
  \IfFileExists{../../../../public/_shared/aiml-sources.tex}{%
    % Source strings for Elements of AIML (used by slides + notes).

\newcommand{\AIMLRepoURL}{https://github.com/tali7c/Elements-of-AIML}

\newcommand{\AIMLLabSources}{%
Provost and Fawcett, \textit{Data Science for Business}.;%
McKinney, \textit{Python for Data Analysis}.;%
WEKA Documentation (Waikato Environment for Knowledge Analysis), accessed 2026-02-28.;%
Matplotlib and Seaborn documentation, accessed 2026-02-28.%
}

\newcommand{\AIMLUnitISources}{%
Rich and Knight, \textit{Artificial Intelligence}, McGraw Hill, 2017.;%
Russell and Norvig, \textit{Artificial Intelligence: A Modern Approach} (4e), Ch. 1--2 (Introduction, Intelligent Agents).;%
Poole and Mackworth, \textit{Artificial Intelligence: Foundations of Computational Agents} (2e), selected sections.%
}

\newcommand{\AIMLUnitIISources}{%
Russell and Norvig, \textit{Artificial Intelligence: A Modern Approach} (4e), Ch. 7--9 (Logical Agents, First-Order Logic, Inference).;%
Huth and Ryan, \textit{Logic in Computer Science} (2e), selected sections on propositional and first-order logic.;%
Genesereth and Nilsson, \textit{Logical Foundations of Artificial Intelligence}, selected chapters.%
}

\newcommand{\AIMLUnitIIISources}{%
Mueller and Massaron, \textit{Machine Learning for Dummies}, For Dummies, 2016.;%
Mitchell, \textit{Machine Learning}, selected chapters on learning basics and evaluation.;%
Scikit-learn documentation, accessed 2026-02-28.%
}

\newcommand{\AIMLUnitIVSources}{%
Mueller and Massaron, \textit{Machine Learning for Dummies}, For Dummies, 2016.;%
James, Witten, Hastie, and Tibshirani, \textit{An Introduction to Statistical Learning} (2e), selected chapters.;%
Scikit-learn documentation, accessed 2026-02-28.%
}

\newcommand{\AIMLUnitVSources}{%
NIST, \textit{AI Risk Management Framework (AI RMF 1.0)}, 2023.;%
Barocas, Hardt, and Narayanan, \textit{Fairness and Machine Learning} (online book), selected sections.;%
Knaflic, \textit{Storytelling with Data}, selected chapters on communicating results.%
}
%
  }{%
    % Source strings for Elements of AIML (used by slides + notes).

\newcommand{\AIMLRepoURL}{https://github.com/tali7c/Elements-of-AIML}

\newcommand{\AIMLLabSources}{%
Provost and Fawcett, \textit{Data Science for Business}.;%
McKinney, \textit{Python for Data Analysis}.;%
WEKA Documentation (Waikato Environment for Knowledge Analysis), accessed 2026-02-28.;%
Matplotlib and Seaborn documentation, accessed 2026-02-28.%
}

\newcommand{\AIMLUnitISources}{%
Rich and Knight, \textit{Artificial Intelligence}, McGraw Hill, 2017.;%
Russell and Norvig, \textit{Artificial Intelligence: A Modern Approach} (4e), Ch. 1--2 (Introduction, Intelligent Agents).;%
Poole and Mackworth, \textit{Artificial Intelligence: Foundations of Computational Agents} (2e), selected sections.%
}

\newcommand{\AIMLUnitIISources}{%
Russell and Norvig, \textit{Artificial Intelligence: A Modern Approach} (4e), Ch. 7--9 (Logical Agents, First-Order Logic, Inference).;%
Huth and Ryan, \textit{Logic in Computer Science} (2e), selected sections on propositional and first-order logic.;%
Genesereth and Nilsson, \textit{Logical Foundations of Artificial Intelligence}, selected chapters.%
}

\newcommand{\AIMLUnitIIISources}{%
Mueller and Massaron, \textit{Machine Learning for Dummies}, For Dummies, 2016.;%
Mitchell, \textit{Machine Learning}, selected chapters on learning basics and evaluation.;%
Scikit-learn documentation, accessed 2026-02-28.%
}

\newcommand{\AIMLUnitIVSources}{%
Mueller and Massaron, \textit{Machine Learning for Dummies}, For Dummies, 2016.;%
James, Witten, Hastie, and Tibshirani, \textit{An Introduction to Statistical Learning} (2e), selected chapters.;%
Scikit-learn documentation, accessed 2026-02-28.%
}

\newcommand{\AIMLUnitVSources}{%
NIST, \textit{AI Risk Management Framework (AI RMF 1.0)}, 2023.;%
Barocas, Hardt, and Narayanan, \textit{Fairness and Machine Learning} (online book), selected sections.;%
Knaflic, \textit{Storytelling with Data}, selected chapters on communicating results.%
}
%
  }%
}


\title{Elements of AIML\\Unit 02 -- Lecture 02 Notes\\Predicate Logic (First-Order Logic)}
\author{Tofik Ali}
\date{\today}

\begin{document}
\maketitle
\tableofcontents

\section{Lecture Overview}
This lecture introduces \textbf{predicate logic} (also called \textbf{first-order logic}, FOL). Compared to propositional logic, FOL can express:
\begin{itemize}[leftmargin=*]
  \item \textbf{objects} (Ali, UPES, Course-1),
  \item \textbf{properties} (Student(x)),
  \item \textbf{relations} between objects (Enrolled(x,y)),
  \item and \textbf{general rules} using quantifiers ($\forall$, $\exists$).
\end{itemize}
\notesources{\AIMLUnitIISources}

\section{How to use these notes (self-study)}
When translating English to FOL, write down (1) your domain, (2) your predicate meanings, and (3) the intended quantifier scope.
Most errors in exams happen because one of these three was implicit in the student's mind but not written in logic.
Try to translate each practice sentence \emph{before} reading the provided answer.

\section{Learning Outcomes}
\begin{itemize}[leftmargin=*]
  \item Explain why propositional logic becomes verbose when there are many objects.
  \item Write basic FOL sentences with predicates, variables, and quantifiers.
  \item Interpret the meaning of universal and existential quantifiers.
  \item Translate simple English statements into FOL and avoid common pitfalls.
\end{itemize}

\section{Keywords}
\begin{itemize}[leftmargin=*]
  \item \textbf{Domain:} the set of objects that variables range over.
  \item \textbf{Constant:} a symbol naming a specific object (e.g., $Ali$).
  \item \textbf{Predicate:} a relation or property returning True/False (e.g., $Student(x)$).
  \item \textbf{Function:} maps objects to objects (e.g., $FatherOf(x)$).
  \item \textbf{Quantifier:} $\forall$ (for all), $\exists$ (there exists).
  \item \textbf{Scope:} the part of the sentence controlled by a quantifier.
\end{itemize}

\section{Abbreviations and symbols}
\begin{itemize}[leftmargin=*]
  \item \textbf{FOL}: First-Order Logic (predicate logic).
  \item \textbf{PL}: Propositional Logic.
  \item $\forall$: universal quantifier (for all).
  \item $\exists$: existential quantifier (there exists).
  \item $\lnot,\land,\lor,\rightarrow$: NOT, AND, OR, IMPLIES.
\end{itemize}

\section{Why predicate logic?}
In propositional logic, we represent facts as atomic symbols. For example:
\begin{itemize}[leftmargin=*]
  \item $p$: ``Ali is a student.''
  \item $q$: ``Sara is a student.''
\end{itemize}
If we have 500 students, we need 500 separate symbols and cannot easily express a rule like:
\begin{quote}
``All students are hardworking.''
\end{quote}
Predicate logic solves this by introducing variables and quantifiers:
\[
  \forall x\; Student(x) \rightarrow WorksHard(x)
\]

\section{Syntax of FOL}
\subsection{Terms}
Terms refer to objects:
\begin{itemize}[leftmargin=*]
  \item constants: $Ali$
  \item variables: $x$
  \item function applications: $FatherOf(x)$
\end{itemize}

\subsection{Predicates vs functions (common confusion)}
\begin{itemize}[leftmargin=*]
  \item A \textbf{predicate} returns True/False (it is a statement).
  \item A \textbf{function} returns an object (it is a term-producing operator).
\end{itemize}
Example: $Parent(Ali,Sara)$ could be a predicate (True/False). $FatherOf(Ali)$ is a function producing an object.

\subsection{Atomic sentences}
An atomic sentence is a predicate applied to terms, e.g.:
\[
  Student(Ali),\quad Enrolled(Ali, CS101),\quad Likes(x, AI)
\]

\subsection{Complex sentences}
Build complex sentences using the same connectives as propositional logic:
\[
  \lnot,\; \land,\; \lor,\; \rightarrow,\; \leftrightarrow
\]
and include quantifiers:
\[
  \forall x\; (Student(x) \rightarrow WorksHard(x))
\]

\subsection{Free vs bound variables (scope)}
\begin{itemize}[leftmargin=*]
  \item A variable is \textbf{bound} if it appears under a quantifier that declares it.
  \item A variable is \textbf{free} if it is not bound by any quantifier.
\end{itemize}
Example: in $\forall x\; Likes(x, AI)$, $x$ is bound. In $Likes(x, AI)$ (alone), $x$ is free and the sentence is not a complete statement unless we specify what $x$ is.

\section{Quantifiers and scope}
\subsection{Universal quantifier ($\forall$)}
$\forall x\; P(x)$ means: for every object $x$ in the domain, $P(x)$ holds.

\subsection{Existential quantifier ($\exists$)}
$\exists x\; P(x)$ means: there exists at least one object $x$ such that $P(x)$ holds.

\subsection{Scope and order}
The \textbf{order} of quantifiers can change meaning.

\paragraph{Example.}
``Every student is enrolled in some course'':
\[
  \forall x\; (Student(x) \rightarrow \exists y\; (Course(y) \land Enrolled(x,y)))
\]
This allows different students to be enrolled in different courses.

``There exists a course that every student is enrolled in'':
\[
  \exists y\; (Course(y) \land \forall x\; (Student(x) \rightarrow Enrolled(x,y)))
\]
This asserts one shared course for all students.

\subsection{Negating quantifiers (very useful later)}
Quantifier negations:
\[
  \lnot(\forall x\; P(x)) \equiv \exists x\; \lnot P(x)
\]
\[
  \lnot(\exists x\; P(x)) \equiv \forall x\; \lnot P(x)
\]
These equivalences are important in clause form conversion and resolution (Unit II Lecture 03).

\section{Semantics: interpreting FOL}
To evaluate truth, we must choose:
\begin{itemize}[leftmargin=*]
  \item a \textbf{domain} (what objects exist),
  \item an \textbf{interpretation} (what constants, functions, and predicates mean).
\end{itemize}
Then sentences can be evaluated as True/False. A \textbf{model} is an interpretation where a set of sentences is true.

\subsection{A tiny semantics example}
Let the domain be the set $\{Ali, Sara\}$ and predicate $Student(\cdot)$.
If the interpretation says:
\[
  Student(Ali)=T,\quad Student(Sara)=F,
\]
then $\exists x\; Student(x)$ is true, while $\forall x\; Student(x)$ is false.

\section{Translation: English $\rightarrow$ FOL}
\subsection{Practical workflow}
\begin{enumerate}[leftmargin=*]
  \item Choose a domain (e.g., people, courses).
  \item Define predicates/functions clearly (e.g., $Student(x)$, $Enrolled(x,y)$).
  \item Identify quantifiers and their scope.
  \item Build the final formula and check if it matches meaning.
\end{enumerate}

\subsection{Common patterns}
\begin{itemize}[leftmargin=*]
  \item ``All A are B'': $\forall x\; A(x) \rightarrow B(x)$
  \item ``Some A is B'': $\exists x\; A(x) \land B(x)$
  \item ``No A is B'': $\forall x\; A(x) \rightarrow \lnot B(x)$
\end{itemize}

\subsection{More patterns (important in practice)}
\begin{itemize}[leftmargin=*]
  \item ``Only A are B'': $\forall x\; B(x) \rightarrow A(x)$
  \item ``At least one A'': $\exists x\; A(x)$
  \item ``At most one A'' (sketch): $\forall x\forall y\; (A(x)\land A(y) \rightarrow x=y)$
  \item ``Exactly one A'' (sketch): (at least one) $\land$ (at most one)
\end{itemize}

\subsection{Practice (with answers)}
Assume domain: all people and courses in a university.

\paragraph{P1. ``All humans are mortal.''}
\textbf{Answer:} $\forall x\; Human(x) \rightarrow Mortal(x)$.

\paragraph{P2. ``Some student likes AI.''}
\textbf{Answer:} $\exists x\; Student(x) \land Likes(x, AI)$.

\paragraph{P3. ``Every student has a mentor who is a faculty member.''}
\textbf{Answer:}
\[
  \forall x\; (Student(x) \rightarrow \exists y\; (Faculty(y) \land MentorOf(y,x)))
\]

\paragraph{P4. ``No student is both absent and present.''}
\textbf{Answer:} $\forall x\; Student(x) \rightarrow \lnot(Absent(x) \land Present(x))$.

\paragraph{P5. ``Only students can access the lab.''}
\textbf{Answer:} $\forall x\; AccessLab(x) \rightarrow Student(x)$.

\paragraph{P6. ``Exactly one student is class representative.''}
\textbf{Answer (one possible encoding):}
\[
  \exists x\; (Student(x)\land Rep(x)) \;\land\;
  \forall y\; ((Student(y)\land Rep(y)) \rightarrow y=x)
\]

\section{Bridge to next lecture: unification idea (preview)}
Resolution in FOL needs \textbf{unification}, which is essentially ``pattern matching'' for terms.
Example:
\[
  Likes(x, AI) \text{ matches } Likes(Ali, AI) \text{ with substitution } \{x/Ali\}.
\]
We will study this formally in Lecture 03.

\section{Practice (with answers)}

\subsection*{P7. Translate and explain scope.}
Translate: ``Every student likes some course.''
\textbf{Answer:} $\forall x\; (Student(x) \rightarrow \exists y\; (Course(y)\land Likes(x,y)))$.
This means each student may like a different course.

\subsection*{P8. Quantifier order changes meaning.}
Translate: ``There exists a course that every student likes.''
\textbf{Answer:} $\exists y\; (Course(y) \land \forall x\; (Student(x)\rightarrow Likes(x,y)))$.
This asserts a single course liked by all students.

\section{Common pitfalls (and how to avoid them)}
\begin{itemize}[leftmargin=*]
  \item \textbf{Wrong implication direction:} ``All A are B'' is $A(x)\rightarrow B(x)$, not $B(x)\rightarrow A(x)$.
  \item \textbf{Missing scope parentheses:} always group the part controlled by a quantifier.
  \item \textbf{Quantifier order confusion:} $\forall x\exists y$ (depends on $x$) vs $\exists y\forall x$ (one $y$ for all).
  \item \textbf{Forgetting the domain:} the truth of a FOL sentence depends on what objects $x$ ranges over.
  \item \textbf{Free variables:} a sentence with a free variable is not a complete statement unless you quantify it.
  \item \textbf{Predicate vs function:} predicates are True/False statements; functions produce objects.
\end{itemize}

\section{Exit questions (with answers)}

\subsection*{Q1. Why is propositional logic insufficient for many knowledge bases?}
\textbf{Answer:} It cannot represent objects and relations compactly; it requires one atom per fact and cannot express general rules using quantifiers.

\subsection*{Q2. Translate: ``All humans are mortal'' and ``Some student likes AI''.}
\textbf{Answer:} $\forall x\; Human(x) \rightarrow Mortal(x)$ and $\exists x\; Student(x) \land Likes(x, AI)$.

\subsection*{Q3. Why does quantifier order matter?}
\textbf{Answer:} $\forall x \exists y$ allows $y$ to depend on $x$; $\exists y \forall x$ requires a single $y$ that works for all $x$. The meaning changes.

\subsection*{Q4. Give one ambiguous English sentence without assumptions.}
\textbf{Answer:} ``Every student solved a problem.'' (same problem for all vs possibly different problems). Without clarifying, there are multiple valid FOL translations.

\section{Summary}
Predicate logic extends propositional logic by adding objects, relations, and quantifiers, enabling compact and expressive knowledge representation.
\notesources{\AIMLUnitIISources}

\end{document}
